\documentclass[a4paper, ngerman, oneside, 10pt]{article}
\usepackage[utf8]{inputenc} % erlaubt Umlaute und Sonderzeichen (in Overleaf standartmäßig aktiv)
\usepackage{graphicx} % um Bilder einzufügen
\usepackage{tabularx} % um Tabellen zu erstellen
\usepackage{amsmath, amsfonts, amssymb} % Mathematische Symbole der American Mathematical Society
\usepackage[a4paper,margin=3cm]{geometry} % um die Seitenränder zu verändern
\usepackage{siunitx} % um Einheiten zu schreiben
\usepackage[ngerman]{babel} % Deutsches Datumsformat, deutsches Hochkomma, etc.
\usepackage{csquotes} % um Anführungszeichen zu setzen
\usepackage{hyperref} % um Referenzen und Links Clickbarmachen
\usepackage{xcolor} % um Farben zu verwenden
    \hypersetup{colorlinks,
        linkcolor={black!50!black},
        citecolor={blue!50!black},
        urlcolor={blue!80!black}}
\usepackage{biblatex} % um Literaturverzeichnisse zu erstellen
    \addbibresource{definitions/literatur.bib} % Pfad zur .bib Datei
\usepackage{eso-pic} % um Bilder im Hintergrund zu platzieren
\usepackage{multicol} % um Text in mehreren Spalten anzuzeigen
\usepackage{float} % um Bilder und Tabellen an einer bestimmten Stelle zu platzieren
\usepackage{booktabs} % um Tabellen zu erstellen
\usepackage{listings} % um Code einzufügen

\usepackage{fancyhdr} % um Kopf- und Fußzeilen zu erstellen
\usepackage{tikz} % um Grafiken zu erstellen
\usepackage{circuitikz} % um Schaltpläne zu erstellen
    \usetikzlibrary{arrows} % um Pfeile in Tikz zu erstellen
%\usepackage{pgfplots} % um Plots zu erstellen
\usepackage[printonlyused, withpage]{acronym} % um Abkürzungen zu erstellen, printonlyused zeigt nur die verwendeten Abkürzungen an, withpage zeigt die Seitenzahl an, auf der die Abkürzung verwendet wird

\pagestyle{fancy}

\def\semester{3}
\def\professor{Dr. Gerda Strutzenberger \& Lukas Neururer \& Dott. Mag. Yunus Schmirander}
\def\gruppe{BA-MGST-23-4D}
\def\datum{\today}
\def\verfasser{Dominic Vogt \& Florian Just \& Fabian Kocher \& Lukas Höpflinger}
\def\titel{Laborbericht 2 - Biosignalverarbeitung}



\begin{document}
\AddToShipoutPictureBG*{%
  \AtPageUpperLeft{%
    \raisebox{-1.3\height}{%
      \includegraphics[width=10cm]{images/MCI_Logo.pdf}%
    }
  }
}

\begin{titlepage}
    \centering
    {\scshape\LARGE Bachelor: Medizin-, Gesundheits- und Sporttechnologie \par}
    \vspace{1cm}
    {\scshape\Large \semester. Semester\par}
    \vspace{1.5cm}
    {\huge\bfseries \titel \par}
    \vspace{0.5cm}
    \vspace{2cm}
    {\Large\bfseries Lehrveranstaltungsleiter: \professor \par}
    \vfill
    {\Large Gruppe: \gruppe\par}
    \vfill
    {\Large Verfasser: \verfasser \par}
    \vspace{1cm}
    {\Large \datum\par}
\end{titlepage}

\newpage
\pagenumbering{Roman}
\setcounter{page}{\numexpr\value{page}+1} % im Leitfaden Inhaltsverzeichnis beginnt mit 2
\addtocontents{toc}{\protect\setcounter{tocdepth}{-1}}
\tableofcontents
\addtocontents{toc}{\protect\setcounter{tocdepth}{3}}

\newpage
\pagenumbering{arabic}

\section{Einleitung}
Dieses Dokument stellt eine Vorlage zur Erstellung von Laborberichten in \LaTeX für den Studiengang Mechatronik dar. Achten Sie darauf, dass Sie die jeweils aktuellste Version verwenden. Richtlinien können sich von Jahr zu Jahr "andern.

\section{Versuchsbeschreibung}
In den meisten F"allen wird es sich anbieten den jeweiligen Versuchsaufbau zu beschreiben. Dies sollte am Anfang jedes Berichtes passieren. Achten Sie darauf, dass falls die wiedergegebenen Bilder nicht die ihrigen sind die Quelle zu zitieren.

\section{Vorausgehende Arbeiten}
Normalerweise sind Sie dazu angehalten sich auf die Laborübung vorzubereiten. Diese Arbeit ist für das Verständnis ihres Berichtes essentiell und sollte ebenfalls der eigentlichen Labordokumentation vorangestellt werden. Im Speziellen betrifft dies
\begin{itemize}
	\item die Messstrategie
	\item allfällige Berechnungen und Abschätzungen
	\item von Ihnen angefertigte Skripte (Quellcode in den Anhang)
	\item vorab getätigte Simulationen
\end{itemize}
Vergessen Sie nicht verwendete Literatur zu zitieren wie z.B.~Skripten, siehe \cite{meh10}.

\section{Beantwortung der Aufgabenstellung}
Erläutern Sie an dieser Stelle \emph{genau} was ihre Aufgabe ist. Am besten durch ein deutlich gekennzeichnetes, wörtliches Zitat:
\begin{quote}
	Kein Gegenstand enthüllt jemals durch die Eigenschaften, die den Sinnen erscheinen, die Ursachen, die ihn hervorgebracht haben, noch die Wirkungen, die aus ihm entspringen werden.
\end{quote}
\rightline{D.~Hume}
Stellen Sie im Anschluss \emph{genau} dar, wie Sie die Aufgabe gelöst haben inklusive Beschreibung der Messanordnung, Dokumentation allfälliger Auffälligkeiten, einer Fehlerabschätzung und ihrem Ergebnis. Ihre Messwerte können auch im Anhang dokumentiert sein. Besonders wichtig ist eine Interpretation des Ergebnisses bzw.~etwaiger auffälliger Abweichungen vom erwarteten Ergebnis.

\section{Grafische Darstellung der Ergebnisse}

Bzgl.~der grafischen Darstellung der Ergebnisse beachten Sie den diesbezüglichen Leitfaden. Hier noch in Kürze die häufigsten Probleme aufgezählt
\begin{itemize}
	\item Jeder Graph muss vom Text referenziert werden. Falls nicht stellt dies seine Notwendigkeit in Frage.
	\item Kurven müssen auch nach einer schwarz-weiß-Kopie unterscheidbar bleiben. Dies geschieht am besten durch Unterschiede im Linienstil (durchgezogen, strichliert, punktiert, $\ldots$).
	\item Werden mehrere Kurven in einen Graphen eingetragen, so ist entweder eine Legende einzufügen (ab 3 Kurven) oder diese in der Bildunterschrift zu unterscheiden (2 Kurven).
	\item Die Achsen müssen beschriftet werden (\textit{Symbol}/Einheit).
	\item Schwarzer Hintergrund ist nicht zulässig!
\end{itemize}

\newpage
\section{Textformatierung}
\textbf{Lorem} ipsum dolor sit amet, consectetur adipiscing elit. Aenean aliquet sapien eu quam blandit, vel auctor eros tincidunt. Proin lobortis, mi et efficitur ultrices, eros lorem ornare justo, sit amet efficitur nibh tortor id purus. Nam semper pharetra quam ullamcorper pretium. Integer non consectetur leo. Nulla viverra dolor vel dolor elementum egestas. Pellentesque habitant morbi senectus et netus et malesuada fames ac turpis egestas. Ut pulvinar maximus interdum. Cras faucibus lorem a nisl porta dapibus. Suspendisse volutpat iaculis magna ut placerat. Praesent sed pellentesque ex. Etiam lacinia elit eu neque imperdiet tristique at et diam. Nullam auctor, erat quis volutpat congue, velit ex finibus est, at lacinia tortor massa non nibh.\break
\textbf{Quisque} interdum efficitur nibh at condimentum. Morbi eget quam blandit, suscipit ipsum id, accumsan arcu. Nam sollicitudin nisl sit amet elit elementum, ac volutpat neque sollicitudin. Sed nec augue enim. Phasellus vel sodales urna. Etiam mattis ultrices metus, eget tincidunt massa rhoncus quis. Vivamus vitae ligula euismod, interdum neque eget, aliquam felis. Duis pharetra consectetur hendrerit. Quisque enim est, volutpat et turpis consectetur, lobortis commodo justo. Integer ipsum ante, commodo sit amet neque ut, elementum vulputate ligula. Sed at massa cursus, pharetra sem sit amet, sagittis nibh. Vivamus orci ex, facilisis porttitor sem eu, commodo eleifend ligula.\\
\textbf{Lorem} ipsum dolor sit amet, consectetur adipiscing elit. Aenean aliquet sapien eu quam blandit, vel auctor eros tincidunt. Proin lobortis, mi et efficitur ultrices, eros lorem ornare justo, sit amet efficitur nibh tortor id purus. Nam semper pharetra quam ullamcorper pretium. Integer non consectetur leo. Nulla viverra dolor vel dolor elementum egestas. Pellentesque habitant morbi tristique senectus et netus et malesuada fames ac turpis egestas. Ut pulvinar maximus interdum. Cras faucibus lorem a nisl porta dapibus. Suspendisse volutpat iaculis magna ut placerat. Praesent sed pellentesque ex. Donec semper pharetra sollicitudin. Etiam lacinia elit eu neque imperdiet tristique at et diam. Nullam auctor, erat quis volutpat congue, velit ex finibus est, at lacinia tortor massa non nibh.\smallbreak
\textbf{Quisque} interdum efficitur nibh at condimentum. Morbi eget quam blandit, suscipit ipsum id, accumsan arcu. Nam sollicitudin nisl sit amet elit elementum, ac volutpat neque sollicitudin. Sed nec augue enim. Phasellus vel sodales urna. Etiam mattis ultrices metus, eget tincidunt massa rhoncus quis. Vivamus vitae ligula euismod, interdum neque eget, aliquam felis. Duis pharetra consectetur hendrerit. Quisque enim est, volutpat et turpis consectetur, lobortis commodo justo. Integer ipsum ante, commodo sit amet neque ut, elementum vulputate ligula. Sed at massa cursus, pharetra sem sit amet, sagittis nibh. Vivamus orci ex, facilisis porttitor sem eu, commodo eleifend ligula.\bigbreak
\textbf{Lorem} ipsum dolor sit amet, consectetur adipiscing elit. Aenean aliquet sapien eu quam blandit, vel auctor eros tincidunt. Proin lobortis, mi et efficitur ultrices, eros lorem ornare justo, sit amet efficitur nibh tortor id purus. Nam semper pharetra quam ullamcorper pretium. Integer non consectetur leo. Nulla viverra dolor vel dolor elementum egestas. Pellentesque habitant morbi tristique senectus et netus et malesuada fames ac turpis egestas. Ut pulvinar maximus interdum. Cras faucibus lorem a nisl porta dapibus. Suspendisse volutpat iaculis magna ut placerat. Praesent sed pellentesque ex. Donec semper pharetra sollicitudin. Etiam lacinia elit eu neque imperdiet tristique at et diam. Nullam auctor, erat quis volutpat congue, velit ex finibus est, at lacinia tortor massa non nibh.

Das ist ein schöner Satz.\\ Das ist ein \emph{schöner} Satz.

\section{sonstige}

\subsection{Text vs. Steuerbefehl}
Das ist ein schöner Satz.\\ Das ist ein \emph{schöner} Satz.

\subsection{Textformatierung genau}
Das \emph{ist \emph{weniger wichtiger} als} andere.\\
Text ist \textbf{fett}.\\
\textit{kursiv}.\\
\textsl{schräg}.\\
\texttt{Sperrdruck}.\\
etc.

\subsection{Umgebungen}
\begin{center}
	Zentrierte Umgebung
\end{center}
\begin{flushleft}
	links
\end{flushleft}
\begin{flushright}
	rechts
\end{flushright}

\subsubsection{Liste}
\begin{itemize}
	\item A
	      \begin{itemize}
		      \item a
		      \item 1
	      \end{itemize}
	\item B
	\item C
\end{itemize}

\subsubsection{Aufzählung}
\begin{enumerate}
	\item A
	\item B
	\item C
\end{enumerate}

\subsubsection{Bilder}
\begin{figure}[H]
	\centering
	\includegraphics[width=0.5\textwidth]{images/MCI_Logo.pdf}
	\caption{Das MCI Logo}
	\label{fig:mci_logo}
\end{figure}


\subsubsection{Tabellen}
% Tipp: https://www.tablesgenerator.com/
\begin{tabular}{lcr}
	1 & 2 & 3 \\
	4 & 5 & 6 \\
	7 & 8 & 9 \\
\end{tabular}

\begin{tabular}{l|c|r}
	1 & 2 & 3 \\
	4 & 5 & 6 \\
	7 & 8 & 9 \\
\end{tabular}

\begin{center}
	\begin{tabular}{|l|c|r|}
		\hline
		1 & 2 & 3 \\
		\hline
		4 & 5 & 6 \\
		\hline
		7 & 8 & 9 \\
		\hline
	\end{tabular}
\end{center}

% Tipp: https://www.tablesgenerator.com/
\begin{table}[H]
	\resizebox{\columnwidth}{!}{%
		\begin{tabular}{@{}llll@{}}
			\textbf{Stakeholder}                                                                                              &
			\textbf{Ziel / Motivation}                                                                                        &
			\textbf{Hürden / Bedenken}                                                                                        &
			\textbf{Nutzen aus eHealth}                                                                                                     \\ \toprule
			\textbf{Herr M.}                                                                                                  &
			\begin{tabular}[c]{@{}l@{}}Selbstständigkeit bewahren,\\ Sicherheit erhöhen\end{tabular}                          &
			\begin{tabular}[c]{@{}l@{}}Ablehnung der Überwachung,\\ Handhabung,\\ Schamgefühle,\\ Technikangst\end{tabular}   &
			\begin{tabular}[c]{@{}l@{}}Frühzeitige Hilfe bei Stürzen,\\ längerer Verbleib zu Hause,\\ Gefühl von Unterstützung\end{tabular} \\ \midrule
			\textbf{Angehörige}                                                                                               &
			\begin{tabular}[c]{@{}l@{}}Entlastung,\\ Sicherheit von Herrn M.\\ gewährleisten\end{tabular}                     &
			\begin{tabular}[c]{@{}l@{}}Misstrauen gegenüber Datenschutz,\\ unnötige Sorgen bei Fehlalarmen\end{tabular}       &
			\begin{tabular}[c]{@{}l@{}}Beruhigung durch Monitoring,\\ aktive Einbindung\\ ohne ständige Präsenz\end{tabular}                \\ \midrule
			\textbf{Hausarzt}                                                                                                 &
			\begin{tabular}[c]{@{}l@{}}Prävention von Komplikationen,\\ Vermeidung von\\ Krankenhausaufenthalten\end{tabular} &
			Alarmmüdigkeit durch Fehlalarme                                                                                   &
			\begin{tabular}[c]{@{}l@{}}Objektive Bewegungsdaten,\\ schnellere Reaktionszeiten,\\ Entscheidungsunterstützung\end{tabular}    \\ \midrule
			\textbf{\acf{rd}}                                                                                                  &
			\begin{tabular}[c]{@{}l@{}}Schnelle und gezielte\\ Versorgung im Notfall\end{tabular}                             &
			\begin{tabular}[c]{@{}l@{}}Fehlalarme,\\ unklare Ortsangaben bei\\ \acs{gps}-Problemen\end{tabular}               &
			\begin{tabular}[c]{@{}l@{}}Frühe Alarmierung,\\ präzise Einsatzkoordination,\\ Zeitgewinn bei echten Notfällen\end{tabular}     \\ \midrule
			\textbf{Krankenkasse}                                                                                             &
			\begin{tabular}[c]{@{}l@{}}Kostenreduktion durch Prävention\\ und kürzere Liegezeiten\end{tabular}                &
			\begin{tabular}[c]{@{}l@{}}Nachweis des Nutzens,\\ Finanzierung von Technik,\\ Datenschutzregelungen\end{tabular} &
			\begin{tabular}[c]{@{}l@{}}Reduktion teurer\\ Krankenhausaufenthalte,\\ langfristige Versorgungseffizienz\end{tabular}
		\end{tabular}%
	}
	\caption{Stakeholder-Perspektiven beim Einsatz eines Sturzerkennungs-Wearables}
	\label{tab:stakeholder}
\end{table}

\clearpage
\section{TikZ und Plots}
\begin{figure}[h!]
	\centering
	\begin{tikzpicture}[>=stealth]
		\node[draw=green, fill=black] (A) at (0,0) {\textcolor{white}{Unser Knoten A}};
		\node[fill=red] (B) at (1,3) {};
		\draw[<->, dashed] (A) -- (B);
	\end{tikzpicture}
	\caption{unsere erste TikZ}
	\label{fig:tiks_1}
\end{figure}

\begin{figure}[h!]
	\centering
	\begin{tikzpicture}[>=stealth]
		\draw[step=1.0, gray, thin] (0,0) grid (5,4);
		\node[circle, fill=red] (A) at (1,2) {};
		\node[fill=yellow] (B) at (2,3) {};
		\draw[<->] (1,1) --++ (2,3);
		\draw[fill=blue, thick] (2,1) rectangle (3,3);
	\end{tikzpicture}
	\caption{grid}
	\label{fig:tiks_2}
\end{figure}

\begin{figure}[h!]
	\centering
	\begin{tikzpicture}[>=stealth]
		\draw[step=1.0, gray, thin] (0,0) grid (5,4);
		\draw[thick,color=blue,domain=0:2,samples=50, dotted] plot (\x, \x^2);
		\foreach \x in {0,1,2,3,4,5}
		\draw (\x,1pt) -- (\x,-1pt) node[anchor=north] {$\x$};
		\foreach \y in {1,...,4}
		\draw (1pt,\y) -- (-1pt,\y) node[anchor=east] {$\y$};

		% \foreach \x in {0,1,2,3,4,5,6,7}
		%     \foreach \y in {1,...,4}
		%         \node (node\x\y) at (\x,\y) {$(\x,\y)$};
	\end{tikzpicture}
	\caption{plot}
	\label{fig:tiks_3}
\end{figure}
\subsection{pgfplots}
\begin{figure}[!ht]
	\centering
	\begin{tikzpicture}
		\begin{axis}[title ={$f(x) =x^2$},
				grid ,
				xlabel ={$x$},
				ylabel ={$f(x)$},]
			\addplot[domain=-3:3,red] {x^2};
		\end{axis}
	\end{tikzpicture}
	\caption{Diese Grafik zeigt die Funktion $f(x) = x^2$}
	\label{fig:pgf1}
\end{figure}

\begin{tikzpicture}
	\begin{axis}[ domain = 0:360,
			samples =80,
			width =10cm ,
			height =7cm,
			xmax =800]
		\addplot[color=red ,mark=*] coordinates {
				(-200,-1)
				(-133,-1.2)
				(-66,-2)
				(0,-2.5)
				(66,-4)
				(133,-5)
				(200,-7)
			};
		\addplot[color=blue,dashed] {sin(x)};
		\addplot[color=green, domain=-200:800] {-4+x/90+ cos(x*2)};
	\end{axis}
\end{tikzpicture}

% Mit echten Daten aus einer .dat Datei
%\begin{tikzpicture}
%	\begin{axis}[
%			title = {Histogramm},
%			ybar,
%			bar width=0.9cm,
%			width =10cm
%		]
%		\addplot[fill=blue!50] file {hist.dat};
%	\end{axis}
%\end{tikzpicture}

\subsubsection{3D}
\begin{tikzpicture}
	\begin{axis}
		\addplot3[surf,domain=0:360,samples=50] {sin(x)*sin(y)};
	\end{axis}
\end{tikzpicture}

\newpage
\section{Multicol}
\begin{multicols}{2}
	Quisque interdum efficitur nibh at condimentum. Morbi eget quam blandit, suscipit ipsum id, accumsan arcu. Nam sollicitudin nisl sit amet elit elementum, ac volutpat neque sollicitudin. Sed nec augue enim. Phasellus vel sodales urna. Etiam mattis ultrices metus, eget tincidunt massa rhoncus quis. Vivamus vitae ligula euismod, interdum neque eget, aliquam felis. Duis pharetra consectetur hendrerit. Quisque enim est, volutpat et turpis consectetur, lobortis commodo justo. Integer ipsum ante, commodo sit amet neque ut, elementum vulputate ligula. Sed at massa cursus, pharetra sem sit amet, sagittis nibh. Vivamus orci ex, facilisis porttitor sem eu, commodo eleifend ligula.
\end{multicols}

\section{Mathematik}
\begin{equation}
	a^{2}+b^{2}=c^{2}
	\label{equ:Pytagoras}
\end{equation}
oder
$a^{2}+b^{2}=c^{2}$ oder
$$a^{2}+b^{2}=c^{2}$$

fortgeschritten
\begin{gather}
	\vec{F}=m\cdot\vec{a} \\
	a^b+b^2=c^2
\end{gather}

\begin{align}
	\vec{F} & =m\cdot\vec{a} \\
	a^b+b^2 & =c^2
\end{align}

\begin{equation}
	\begin{aligned}
		\vec{F} & =m\cdot\vec{a} \\
		a^b+b^2 & =c^2
	\end{aligned}
\end{equation}

\begin{equation}
	\begin{gathered}
		\vec{F}=m\cdot\vec{a} \\
		a^b+b^2=c^2
	\end{gathered}
\end{equation}

\begin{subequations}
	\begin{align}
		\vec{F} & =m\cdot\vec{a} \\
		a^b+b^2 & =c^2
	\end{align}
\end{subequations}

\begin{equation}
	1+2+3+a+b+c+d+e+f+g+h+i+j+k+l=m+n+o+p+q+r+s+t+u+v+w+x+y+z+1+2+3
\end{equation}

\begin{multline}
	1+2+3+a+b+c+d+e+f+g+h+i+j+k+l= \\
	m+n+o+p+q+r+s+t+u+v+w+x+y+z+1+2+3
\end{multline}

\begin{equation}
	\vec{v}=
	\begin{pmatrix}
		a \\ b \\ -c
	\end{pmatrix}
\end{equation}

\begin{equation}
	\vec{A}=
	\begin{pmatrix}
		a & d \\
		b & e \\
		c & f
	\end{pmatrix}
\end{equation}

\begin{equation}
	\begin{array}[t]{c c c c c}
		a^2      & + & b^2      & = & c^2 \\
		\sin^2 x & + & \cos^2 x & = & 1
	\end{array}
\end{equation}


\begin{equation}
	\int f(x)\,\text{d}x
\end{equation}

\begin{equation}
	\dv{x} \\
	\dv[3]{x} \\
	\pdv{z}
\end{equation}

\begin{equation}
	\frac{a}{b+\frac{a}{b+\frac{a}{b}}}
\end{equation}

\begin{equation}
	\dfrac{a}{b+\dfrac{a}{b+\dfrac{a}{b}}}
\end{equation}

\begin{equation}
	\boxed{
	\underbrace{y_{Res}(x,t)}_{\text{Auslenkung}} =
	\underbrace{2A \cdot \cos \left( k \cdot x \right)}_{\text{Amplitude am Ort } x} \cdot
	\underbrace{\cos \left( \omega \cdot t + \frac{\varphi}{4} \right)}_{\text{Schwingungs-term}}
	}
\end{equation}

\begin{equation}
	\begin{array}[t]{c c c c c c c}
		3x_1  & - & 3x_2  & + & 3x_3 & = & 0 \\
		8x_1  & + & 10x_2 & + & 2x_3 & = & 6 \\
		-2x_1 & + & x_2   & - & 3x_3 & = & 5
	\end{array}
\end{equation}

\begin{equation}
	\begin{pmatrix}
		3 & 2 & 1 \\
		1 & 0 & 2
	\end{pmatrix}
	\cdot
	\begin{pmatrix}
		1 \\ 0 \\ 4
	\end{pmatrix}
	=
	\begin{pmatrix}
		7 \\ 9
	\end{pmatrix}
\end{equation}

\subsection{SI-Einheiten}
\textbf{Lithium-Polymer-Batterie (\boldmath$\SI{1200}{\milli\ampere{}\hour}$, $\SI{3,7}{\volt}$):}

\begin{itemize}
	\item \num{13,2} oder \num{13.2}
	\item \num{3,78e-5}
	\item \num{3,5 \pm 0,2}
	\item \numlist{1,5;2,4;7,0}
	\item \numrange{1e5}{7e6}
	\item \ang{45;7;22}
\end{itemize}

\begin{itemize}
	\item \si{\newton}
	\item \si{\meter} oder \si{\m}
	\item \si{\m \per \s}
	\item \si{\m \per \s \squared}
	\item \si{\m \per \s \tothe{6}}
\end{itemize}

\si{\seppi}

\begin{itemize}
	\item $F=\SI{25,3 \pm 0,8}{\newton}$
	\item $v=\SI{2e10}{\m \per \s}$
	\item \SIlist{2,5;5,4;6,7}{\watt}
	\item \SIrange{3,5e-3}{2}{\kilo\g}
\end{itemize}

\subsection{package:mchem}
\begin{itemize}
	\item \ce{Ag}
	\item \ce{H2O} \ce{H2SO4}
	\item \ce{Ag^2+}
	\item \ce{Fe^{II}}
	\item \ce{^4_2He} \ce{^{14}C}
	\item \ce{A\bond{=}B}
	\item \ce{A -> B} \ce{A <--> B}
\end{itemize}


\section{Referenzieren}\label{sec:ref}
Gleichung \ref{equ:Pytagoras} ist auf Seite \pageref{equ:Pytagoras}. Jetzt sind wir auf Seite \pageref{sec:ref}\\
\href{mailto:mail@mail.mail}{Email@Email.com}\\
Test \cite{Kleiner.2022}\\
\href{https://www.mci.edu/de/}{MCI}\\
\href{mailto:jkdshf@dlsijkj.at}{E-mail}

\qrcode[hyperlink, height=2cm]{https://www.mci.edu/de/}

\subsection{Abkürzungen}

\begin{itemize}
	\item \ac{led}		% Langform mit (Abkürzung) beim ersten mal
	\item \ac{led}		% nur Abkürzung beim zweiten mal
	\item \acf{led}		% immer Langform mit (Abkürzung)
	\item \acl{led}		% immer nur Langform
	\item \acs{led}		% immer nur Kurzform
	\item \acp{led}		% Pluralform der Abkürzung
	\item \aclp{led}	% Pluralform der Langform
	\item \acfp{led}	% Pluralform der Langform mit (Abkürzung)
	\item \acp{rd}		% selbst definierte Pluralfoem
\end{itemize}





% Platz für Verzeichnisse
\clearpage
\pagenumbering{Roman}
\setcounter{page}{\numexpr\value{RomanPagenumber}}
\listoffigures
\addcontentsline{toc}{section}{Abbildungsverzeichnis}
\newpage
\listoftables
\addcontentsline{toc}{section}{Tabellenverzeichnis}
\newpage
\printbibliography[title={Literaturverzeichnis}]
\addcontentsline{toc}{section}{Literaturverzeichnis}
\newpage
\section*{Abkürzungsverzeichnis}
\addcontentsline{toc}{section}{Abkürzungsverzeichnis}
\begin{acronym}[DSGVO] % Longest Acronym here!
	\acro{dsgvo}[DSGVO]{Datenschutz-Grundverordnung}
	\acro{mci}[MCI]{Management Center Innsbruck}
	\acro{mgst}[MGST]{Medizin-, Gesundheits- und Sporttechnologie}
	\acro{bsc}[BSc]{Bachelor of Science}
 	\acro{msc}[MSc]{Master of Science}
 	\acro{phd}[PhD]{Doctor of Philosophy}
	\acro{ing}[Ing]{Ingenieur}
	\acro{adc}[ADC]{Analog-Digital Converter}
	\acro{led}[LED]{Light Emitting Diode}
	\acro{usb}[USB]{Universal Serial Bus}
	\acro{pwm}[PWM]{Pulse Width Modulation}
	\acro{cad}[CAD]{Computer-Aided Drawing}
	\acro{gps}[GPS]{Global Positioning System}
	\acro{ekg}[EKG]{Elektrokardiogramm}
		\acroplural{ekg}[EKGs]{Elektrokardiogramme} %Definition der Pluralform wenn nötig
\end{acronym}
\end{document}
